\chapter{字符串命令解析器设计} \label{ch:command-parser}

人机交互不只有按键和菜单。当设备需要接收 \texttt{"led on"}、\texttt{"add 12 30"} 这类文本命令时，就需要一个字符串命令解析器：它把一行字符串拆成“命令名 + 参数”，再调用对应的处理函数。本章先介绍最直接的字符串比较方法，再引出词法分析，最后用哈希表构建可扩展的命令注册表。在此基础上，本章还会把参数从“字符串数组”升级为带有动态类型的 \texttt{Param\_T}，让命令处理函数可以直接用 \texttt{INT\_PARAM}、\texttt{FLOAT\_PARAM} 和 \texttt{STR\_PARAM} 读取参数。完整代码位于 \texttt{Chapter-command-parser/字符串命令解析器/} 目录。

\medskip \textbf{本章涉及到的实验内容：}
\begin{enumerate}
\item 字符串比较命令分发；
\item 空格或制表符参数切分（词法分析）；
\item 基于哈希表 + 线性探测的命令注册与查找。
\item 动态类型解析（INT、UINT、FLOAT 和 STRING）；
\item 通过适配器调用 \texttt{LED\_CMD(int)}、\texttt{Car\_Run(float, float, float)} 等不同签名的业务函数。
\end{enumerate}

% ---------------------------------------------------------------------
\section{常规解析方法}

最简单的方法是把整行命令与预定义字符串逐个比较：

\begin{lstlisting}[style=c]
if (strcmp(line, "led on") == 0) {
    Led_On();
} else if (strcmp(line, "led off") == 0) {
    Led_Off();
} else if (strcmp(line, "help") == 0) {
    PrintHelp();
} else {
    printf("unknown command\n");
}
\end{lstlisting}

这种写法在小规模演示中很直观，但存在三个问题：

\begin{enumerate}
\item \textbf{没有参数概念}：\texttt{"led on"} 与 \texttt{"led off"} 是两条完全独立的命令，命令一多就会成倍膨胀；
\item \textbf{代码无法扩展}：新增一条命令必须修改主分发函数，违反“新增不改旧”的模块化原则；
\item \textbf{性能线性}：命令数量达到几十条后，每次都要从头做字符串比较。
\end{enumerate}

\begin{table}[htbp]
\centering \caption{字符串命令解析的三种常用方式对比} \label{tab:cmd-method}
    \begin{tabular}{lllll}
        \toprule
        \textbf{方法} & \textbf{参数切分} & \textbf{查找方式} & \textbf{新增命令} & \textbf{适合规模} \\
        \midrule
        if-else 字符串比较 & 无 & 顺序 strcmp & 修改主函数 & 少量固定命令 \\
        strcmp 命令表 & 无 & 线性遍历表 & 追加表项 & 中等规模 \\
        词法分析 + 哈希表 & 有 & 哈希查找 & 注册表项 & 命令较多且频繁 \\
        \bottomrule
    \end{tabular}
\end{table}

因此工程化做法是“命令表驱动”：命令名 + 处理函数指针构成一张表，解析器只负责查表调用。本章示例中，命令表定义如下：

\begin{lstlisting}[style=c]
typedef int8_t (*CmdHandler_F)(uint8_t argc, Param_T *argv);

typedef struct __CmdItem_T
{
    char name[CMD_NAME_MAX]; /* 命令名 */
    CmdHandler_F handler;    /* 处理函数 */
    uint8_t used;            /* 是否已注册 */
} CmdItem_T;
\end{lstlisting}

% ---------------------------------------------------------------------
\section{词法分析器}

“词法分析”在这里指把一行命令切分成最小有意义的单元（Token）。对于空格分隔的命令行，只需两步：

\begin{enumerate}
\item 跳过连续空格或制表符；
\item 记录当前参数起点，一直走到下一个空白或行尾，把该位置替换为 \texttt{'\textbackslash0'}。
\end{enumerate}

示例实现如下（\texttt{argv} 指向原缓冲区中的各个参数起点）：

\begin{lstlisting}[style=c]
static uint8_t Cmd_Tokenize(char *pLine, char *argv[], uint8_t maxArgs)
{
    uint8_t argc = 0;
    char *p = pLine;

    while (*p != '\0') {
        while ((*p == ' ') || (*p == '\t')) {
            p++;                 /* 跳过空白 */
        }
        if (*p == '\0') {
            break;
        }
        if (argc >= maxArgs) {
            break;               /* 防止 argv 越界 */
        }

        argv[argc++] = p;        /* 记录参数起点 */

        while ((*p != '\0') && (*p != ' ') && (*p != '\t')) {
            p++;                 /* 移动到参数末尾 */
        }
        if (*p != '\0') {
            *p = '\0';           /* 用结束符切开参数 */
            p++;
        }
    }
    return argc;
}
\end{lstlisting}

切分后，\texttt{"led on"} 会变成 \texttt{argv[0] = "led"}、\texttt{argv[1] = "on"}。词法分析器的要点是：\textbf{边界检查比算法本身更重要}——行太长、参数太多、空命令都要在这一层拦截，避免处理函数访问越界。

\begin{figure}[htbp]
\centering
    \begin{tikzpicture}[
        tok/.style={draw=structurecolor, thick, rounded corners=2pt,
                    minimum width=2.2cm, minimum height=0.9cm, align=center,
                    fill=structurecolor!8, font=\small},
        lab/.style={font=\small},
        arrow/.style={-stealth, thick}
    ]
        \node[tok] (input) at (0,1.8) {命令字符串\\ \texttt{"led on"}};

        \node[tok] (t0) at (-2.2,-1.2) {\texttt{argv[0]}\\ \texttt{"led"}};
        \node[tok] (t1) at (2.2,-1.2)  {\texttt{argv[1]}\\ \texttt{"on"}};

        \draw[arrow] (input.south) -- node[left]  {} (t0.north);
        \draw[arrow] (input.south) -- node[right] {} (t1.north);

        \node[lab] at (0,-2.7) {\texttt{argc = 2}，每个参数以 \texttt{'\textbackslash0'} 结尾};
    \end{tikzpicture}
\caption{词法分析：一行命令被切分成命令名与参数} \label{fig:cmd-token}
\end{figure}

% ---------------------------------------------------------------------
\section{HASH 表的构建}

当命令数量较多时，顺序查找 \texttt{"led"}、\texttt{"add"}、\texttt{"help"}……效率低。可以让命令名先经过一个哈希函数映射到数组下标，再在对应位置插入与查找。哈希冲突用\textbf{线性探测}解决：

\begin{lstlisting}[style=c]
static uint8_t Cmd_CalcHash(const char *pName)
{
    uint32_t hash = 2166136261u;   /* FNV-1a 初始值 */
    const uint8_t *p = (const uint8_t *)pName;

    while (*p != '\0') {
        hash ^= (uint32_t)*p++;
        hash *= 16777619u;
    }
    return (uint8_t)(hash % CMD_TABLE_CAP);
}
\end{lstlisting}

两个不同命令名可能算出同一个下标（冲突），线性探测的处理方式是：从哈希位置开始向后找第一个空槽，把命令放进去；查找时也从哈希位置开始，遇到相同命令名返回，遇到空槽说明命令不存在：

\begin{lstlisting}[style=c]
/* 注册：找到空槽位或同名项 */
start = Cmd_CalcHash(pCmdName);
index = start;
while (loop < CMD_TABLE_CAP) {
    if (table[index].used == 0u) {
        /* 插入新命令 */
        break;
    }
    if (strcmp(table[index].name, pCmdName) == 0) {
        /* 重复注册，更新 handler */
        break;
    }
    index = (index + 1u) % CMD_TABLE_CAP;   /* 线性探测 */
    loop++;
}
\end{lstlisting}

\begin{figure}[htbp]
\centering
    \begin{tikzpicture}[
        cell/.style={draw=structurecolor, thick, minimum width=1.5cm,
                     minimum height=0.9cm, align=center, font=\small},
        used/.style={fill=structurecolor!20},
        lab/.style={font=\small},
        arrow/.style={-stealth, thick}
    ]
        \node[cell]       (c0) at (0,0)   {0};
        \node[cell]       (c1) at (1.6,0) {1};
        \node[cell, used] (c2) at (3.2,0) {help};
        \node[cell]       (c3) at (4.8,0) {3};
        \node[cell, used] (c4) at (6.4,0) {add};
        \node[cell, used] (c5) at (8.0,0) {led};
        \node[cell]       (c6) at (9.6,0) {6};
        \node[cell]       (c7) at (11.2,0){7};

        \node[lab] at (6.4,1.2) {\texttt{hash("add") = 4}};
        \node[lab] at (8.0,-1.2) {\texttt{hash("led") = 4，探测到 5}};
        \node[lab, red] at (8.5,1.2) {冲突};

        \draw[arrow, red] (8.15,0.9) -- (8.95,0.9);
    \end{tikzpicture}
\caption{哈希表线性探测：命令冲突时向后寻找空槽} \label{fig:hash-linear}
\end{figure}

线性探测适合“启动时注册、运行期只查找”的场景：命令表在初始化阶段填好，之后不再删除，因此不需要墓碑（tombstone）标记。哈希表的优点是把“字符串比较”的次数从 $O(n)$ 降为接近 $O(1)$，缺点是必须保证表容量大于实际命令数，否则会退化为线性查找。

% ---------------------------------------------------------------------
\section{命令注册与执行流程}

综合以上两部分，命令解析器的公共 API 只需要三个函数：

\begin{enumerate}
\item \texttt{Cmd\_ParserInit}：清空命令表；
\item \texttt{Cmd\_ParserRegister}：注册命令名与处理函数；
\item \texttt{Cmd\_ParserRun}：接收一行字符串，切词后查表并调用处理函数。
\end{enumerate}

\texttt{Cmd\_ParserRun} 的执行流程：

\begin{enumerate}
\item 检查行长度，复制到内部缓冲（不能直接修改 \texttt{const} 输入）；
\item 调用 \texttt{Cmd\_Tokenize} 得到 \texttt{argc} 与 \texttt{argv}；
\item 用 \texttt{argv[0]} 查哈希表；
\item 找到后调用 \texttt{handler(argc, argv)}，并透传其返回码。
\end{enumerate}

\begin{figure}[htbp]
\centering
    \begin{tikzpicture}[
        box/.style={draw=structurecolor, thick, rounded corners=2pt,
                    minimum width=3.6cm, minimum height=0.9cm, align=center,
                    fill=structurecolor!8, font=\small},
        arrow/.style={-stealth, thick}
    ]
        \node[box] (in)  at (0,3.0)   {命令行字符串};
        \node[box] (len) at (0,1.5)   {长度检查与复制};
        \node[box] (tok) at (0,0)     {\texttt{Cmd\_Tokenize}};
        \node[box] (hash) at (0,-1.5) {哈希查表 \texttt{argv[0]}};
        \node[box] (run) at (0,-3.0)  {\texttt{handler(argc, argv)}};

        \draw[arrow] (in.south)   -- (len.north);
        \draw[arrow] (len.south)  -- (tok.north);
        \draw[arrow] (tok.south)  -- (hash.north);
        \draw[arrow] (hash.south) -- (run.north);
    \end{tikzpicture}
\caption{\texttt{Cmd\_ParserRun} 的处理流程} \label{fig:cmd-flow}
\end{figure}

处理函数示例：

\begin{lstlisting}[style=c]
static int8_t Cmd_HandleLed(uint8_t argc, Param_T *argv)
{
    if (argc < 2u) {
        printf("usage: led on|off\n");
        return CMD_ERR_EXEC;
    }
    if (strcmp(STR_PARAM(argv, 1), "on") == 0) {
        Led_On();
        return CMD_ERR_OK;
    }
    if (strcmp(STR_PARAM(argv, 1), "off") == 0) {
        Led_Off();
        return CMD_ERR_OK;
    }
    return CMD_ERR_EXEC;
}
\end{lstlisting}

% ---------------------------------------------------------------------
\section{第二部分：动态类型参数解析（Param\_T）}

如果 handler 只能收到字符串，那么每个命令都要自己写 \texttt{atoi} 与 \texttt{atof}，而且很容易漏掉格式错误。更好的做法是让解析器在切词后立刻尝试识别每个参数的类型，把“参数”封装成 \texttt{Param\_T}：

\begin{lstlisting}[style=c]
typedef enum __ParamType_T
{
    PARAM_TYPE_NONE = 0,
    PARAM_TYPE_INT,       /* 12、-3 */
    PARAM_TYPE_UINT,      /* 0x10、255 */
    PARAM_TYPE_FLOAT,     /* 3.5、1e2 */
    PARAM_TYPE_STRING     /* led、on */
} ParamType_T;

typedef struct __Param_T
{
    ParamType_T type;              /* 识别出的参数类型 */
    char text[CMD_PARAM_TEXT_MAX]; /* 原始字符串 */
    int32_t intValue;              /* 按整数解释 */
    uint32_t uintValue;            /* 按无符号整数解释 */
    float floatValue;              /* 按浮点数解释 */
} Param_T;
\end{lstlisting}

之所以同时保存三种数值而不是只保存一种，是因为不同命令对同一文本的解释不同：对于 \texttt{"3.9"}，\texttt{INT\_PARAM} 希望得到 3，\texttt{FLOAT\_PARAM} 希望得到 3.9；对于 \texttt{"12"}，\texttt{FLOAT\_PARAM} 又希望得到 12.0。三种解释字段独立存放后，handler 想按什么类型读都可以。

参数识别规则如下：

\begin{table}[htbp]
\centering \caption{动态类型识别规则示例} \label{tab:param-type}
    \begin{tabular}{llll}
        \toprule
        \textbf{原始文本} & \textbf{type} & \textbf{INT\_PARAM} & \textbf{FLOAT\_PARAM} \\
        \midrule
        \texttt{12}      & INT   & 12        & 12.0 \\
        \texttt{-3}      & INT   & -3        & -3.0 \\
        \texttt{0x10}    & UINT  & 16        & 16.0 \\
        \texttt{3.9}     & FLOAT & 3         & 3.9 \\
        \texttt{1e2}     & FLOAT & 100       & 100.0 \\
        \texttt{led}     & STRING & 无效      & 无效 \\
        \bottomrule
    \end{tabular}
\end{table}

处理函数中读取参数只需使用宏：

\begin{lstlisting}[style=c]
int   n     = INT_PARAM(argv, 1);    /* 以 int 类型读取 argv[1] */
uint32_t u  = UINT_PARAM(argv, 2);   /* 以无符号整数读取 argv[2] */
float speed = FLOAT_PARAM(argv, 3);  /* 以 float 读取 argv[3] */
const char *s = STR_PARAM(argv, 4);  /* 以字符串读取 argv[4] */
\end{lstlisting}

\begin{anybox}{宏的参数名可以任意取}
写成 \texttt{INT\_PARAM(argv, 1)} 或 \texttt{INT\_PARAM(paramList, 1)} 效果相同，只要第一个参数是 \texttt{Param\_T} 数组即可。
\end{anybox}

\begin{anybox}{动态类型解释}
\texttt{Param\_T} 只是把一段文本同时解析成整数、浮点与字符串三种候选值；最终按哪种类型读取，由命令处理函数选择对应的 \texttt{INT\_PARAM}、\texttt{FLOAT\_PARAM} 或 \texttt{STR\_PARAM} 决定。
\end{anybox}
% ---------------------------------------------------------------------
\section{第三部分：适配不同的函数列表}

有了 \texttt{Param\_T}，命令层与业务层就可以分离了。假设设备层已经有如下两个业务函数，它们的签名完全不同：

\begin{lstlisting}[style=c]
/* 业务函数 1：只接收一个 int */
static void Led_Cmd(int n)
{
    printf("LED_CMD(%d)\n", n);
}

/* 业务函数 2：接收三个 float */
static void Car_Run(float x, float y, float speed)
{
    printf("Car_Run(x=%.2f, y=%.2f, speed=%.2f)\n", x, y, speed);
}
\end{lstlisting}

C 语言在运行时并不知道命令表里的 handler 应该调用哪个函数、参数是什么类型，因此每个命令需要写一个“适配器”：负责检查参数个数，把 \texttt{Param\_T} 转成具体函数需要的实参。

\begin{lstlisting}[style=c]
/* 适配器 1：intcmd 7 -> Led_Cmd(7) */
static int8_t Cmd_HandleIntCmd(uint8_t argc, Param_T *argv)
{
    if (argc < 2u) {
        printf("usage: intcmd n\n");
        return CMD_ERR_EXEC;
    }
    Led_Cmd(INT_PARAM(argv, 1));
    return CMD_ERR_OK;
}

/* 适配器 2：car x y speed -> Car_Run(x, y, speed) */
static int8_t Cmd_HandleCar(uint8_t argc, Param_T *argv)
{
    if (argc < 4u) {
        printf("usage: car x y speed\n");
        return CMD_ERR_EXEC;
    }
    Car_Run(FLOAT_PARAM(argv, 1),
            FLOAT_PARAM(argv, 2),
            FLOAT_PARAM(argv, 3));
    return CMD_ERR_OK;
}
\end{lstlisting}

这里的“函数列表”就是命令注册表。每接入一个业务函数，只需要追加一个适配器并注册一次：

\begin{lstlisting}[style=c]
Cmd_ParserInit(&parser);

Cmd_ParserRegister(&parser, "led",    Cmd_HandleLed);      /* 字符串参数 */
Cmd_ParserRegister(&parser, "intcmd", Cmd_HandleIntCmd);   /* int 参数 */
Cmd_ParserRegister(&parser, "car",    Cmd_HandleCar);      /* 3 个 float 参数 */
Cmd_ParserRegister(&parser, "add",    Cmd_HandleAdd);      /* 2 个 int 参数 */
\end{lstlisting}

注册表与业务函数的对应关系如表~\ref{tab:cmd-adapter}~所示。

\begin{table}[htbp]
\centering \caption{命令 -> 适配器 -> 业务函数} \label{tab:cmd-adapter}
    \begin{tabular}{llll}
        \toprule
        \textbf{命令示例} & \textbf{适配器} & \textbf{业务函数} & \textbf{参数类型} \\
        \midrule
        \texttt{led on} & \texttt{Cmd\_HandleLed}   & \texttt{Led\_On()} & STRING \\
        \texttt{intcmd 7} & \texttt{Cmd\_HandleIntCmd} & \texttt{Led\_Cmd(int)} & INT \\
        \texttt{car 1.5 2.5 0.8} & \texttt{Cmd\_HandleCar} & \texttt{Car\_Run(float,float,float)} & FLOAT \\
        \texttt{add 12 30} & \texttt{Cmd\_HandleAdd} & 加法实现 & INT \\
        \bottomrule
    \end{tabular}
\end{table}

需要说明的是：\textbf{在理论上和部分工程实践中，确实存在两种“无适配器”的框架方案}，它们试图让命令解析器直接调用任意签名的业务函数，而无需手写适配器。

\subsection*{方案一：参数描述表 + 通用调度器（重型框架）}

在命令表中额外记录业务函数的参数类型序列（如字符串 \texttt{"iff"} 表示 \texttt{int, float, float}），调度器解析完参数后，按类型序列从 \texttt{Param\_T} 数组中逐一提取数值，再通过函数指针调用业务函数——本质上是一个简化版的动态函数绑定。一些用于 PC 或高性能嵌入式的脚本引擎（如 Lua、MicroPython 的 C 绑定层）即采用此类思路。

示意代码如下：

\begin{lstlisting}[style=c]
/* 参数描述表节点 */
typedef struct {
	const char *cmdName;
	const char *typeDesc;      /* 如 "iff" 表示 int, float, float */
	void (*func)(void);        /* 业务函数指针（泛型） */
} CmdDesc_T;

/* 通用调度器（示意） */
static int8_t Cmd_DispatchGeneric(CmdDesc_T *desc, Param_T *args, uint8_t argc)
{
	/* 按 typeDesc 逐个提取参数，此处仅为示意，实际需处理类型匹配与错误 */
	for (int i = 0; i < strlen(desc->typeDesc); i++) {
		char t = desc->typeDesc[i];
		switch (t) {
			case 'i': /* 提取 INT_PARAM(args, i) */ break;
			case 'f': /* 提取 FLOAT_PARAM(args, i) */ break;
			case 's': /* 提取 STR_PARAM(args, i) */ break;
		}
	}
	/* 调用 desc->func(...) —— 但 C 语言无法在运行时构造任意签名的调用，
	实际需借助宏展开或汇编胶水，此处省略具体实现细节。 */
	return CMD_ERR_OK;
}
\end{lstlisting}

但该方案在\textbf{裸机嵌入式环境}中有三个显著代价：

\begin{enumerate}
	\item \textbf{ROM 开销显著增加}：每个命令除了名字和处理函数指针外，还要存储参数个数、类型序列、可能的默认值等描述信息，命令越多，表的体积越大；
	\item \textbf{运行时开销不可忽视}：通用调度器必须根据类型序列动态决定如何从 \texttt{Param\_T} 提取数据，并构造函数调用栈（或搬运参数到寄存器），每次调用的指令周期远多于手写适配器的几条赋值语句；
	\item \textbf{调试复杂度飙升}：通用调度器内部充斥着类型分发、循环与强制转换，一旦参数数量或类型匹配错误，问题往往在运行时才暴露，排查难度远大于单步调试一个显式适配器函数。
\end{enumerate}

\subsection*{方案二：函数指针强制转换（利用 ABI 特性）}

另一种更“激进”的做法是放弃参数描述表，直接把业务函数指针强制转换为形参全为 \texttt{int32\_t} 的通用类型（如 \texttt{void (*)(int32\_t, int32\_t, int32\_t)}），调用时把 \texttt{float} 或 \texttt{int} 都当作 \texttt{int32\_t} 的位模式传入。

\begin{lstlisting}[style=c]
/* 假设业务函数 Car_Run(float x, float y, float speed) */
void (*p)(int32_t, int32_t, int32_t) = (void*)Car_Run;

/* 提取 float 的位模式当作 int32_t 传递 */
int32_t a0 = *(int32_t*)&FLOAT_VAL(args[0]);
int32_t a1 = *(int32_t*)&FLOAT_VAL(args[1]);
int32_t a2 = *(int32_t*)&FLOAT_VAL(args[2]);

p(a0, a1, a2);   /* 在 soft-float ABI 下能“蒙混过关” */
\end{lstlisting}

该方法在 \textbf{soft-float ABI 的 ARM 平台} 上能正常执行，因为浮点数和整数通过相同的通用寄存器传递，函数体内的机器码会正确解释这些位模式。

但是，这种做法存在严重缺陷：

\begin{enumerate}
	\item \textbf{平台依赖性极强}：一旦启用硬浮点 ABI（\texttt{-mfloat-abi=hard}），浮点数通过专用 FPU 寄存器传递，而整数通过通用寄存器，此时强制调用会传入垃圾值，程序立即崩溃；在 x86/x64 等平台更是完全不可行；
	\item \textbf{C 语言未定义行为}：调用一个与函数指针类型不兼容的函数属于标准未定义行为，编译器优化（如严格别名、内联）可能产生完全意料之外的代码；
	\item \textbf{可维护性极差}：阅读代码的人无法从函数指针类型得知实际调用签名，调试器无法正确显示参数值，且难以进行静态检查。
\end{enumerate}

因此，这种技巧仅在高度确定（软浮点、已知编译器）且对代码体积有极致要求的特殊场合偶有使用，\textbf{绝不能作为通用方案推荐}。

\subsection*{本书的选择}

鉴于此，本章面向资源受限且平台多样的裸机环境，\textbf{不引入任何“无适配器”框架}，而是坚持“每个命令一个显式适配器”的轻量级方案。对于 10~30 条命令的典型规模，适配器的额外代码量完全可控（每条约 5~10 行），且每个适配器就是一个普通函数，断点单步、参数边界检查都一目了然。新增业务函数时，照着既有模板复制粘贴并修改参数读取宏即可，无需引入复杂的元编程工具或依赖底层 ABI 细节。

用适配器作为“显式胶水”来换取代码可读性、可移植性与调试便利性，是嵌入式工程实践中务实的取舍。

% ---------------------------------------------------------------------
\section{完整示例代码}

完整可运行的代码位于 \texttt{Chapter-command-parser/字符串命令解析器/}：

\begin{itemize}
\item \texttt{cmd\_parser.h}：命令表、\texttt{Param\_T}、错误码、API；
\item \texttt{cmd\_parser.c}：哈希函数、动态类型切分、命令注册与查找；
\item \texttt{main.c}：注册 \texttt{help}、\texttt{led}、\texttt{intcmd}、\texttt{car}、\texttt{add} 和 \texttt{types} 并测试。
\end{itemize}

在 PC 上使用 \texttt{gcc cmd\_parser.c main.c} 即可编译运行，输出结果如下：

\begin{lstlisting}
>>> help
>>> led on
LED 已打开
>>> intcmd 7
LED_CMD(7)
>>> car 1.5 2.5 0.8
Car_Run(x=1.50, y=2.50, speed=0.80)
>>> add 12 30
12 + 30 = 42
>>> types 12 -3 0x10 3.5 led
argv[1] = 12       type = INT
argv[2] = -3       type = INT
argv[3] = 0x10     type = UINT
argv[4] = 3.5      type = FLOAT
argv[5] = led      type = STRING
>>> unknown_cmd
未知命令
\end{lstlisting}

把处理函数中的 \texttt{printf} 换成真实外设操作后，这套代码就能直接移植到串口调试命令行中。

% ---------------------------------------------------------------------
\section{单片机实验：串口命令解析移植与实测}

本节实验把 PC 上的命令解析器接到单片机串口，通过电脑串口助手或调试终端输入命令。实验代码参考 \texttt{Chapter-command-parser/字符串命令解析器/}，作者先给出实验框架，后续补充接线照片与运行日志。

\begin{anybox}{实验目的}

\begin{enumerate}
\item 在单片机串口中断或 FIFO 之上实现“收到一行 -> 解析 -> 执行”；
\item 验证 \texttt{INT\_PARAM}、\texttt{FLOAT\_PARAM} 和 \texttt{STR\_PARAM} 对真实参数的解析；
\item 通过命令直接控制 LED、小车等业务函数。
\end{enumerate}

\end{anybox}
\begin{anybox}{建议实验步骤}

\begin{enumerate}
\item 把命令解析源码加入工程，接入串口 FIFO 接收；
\item 注册 \texttt{led}、\texttt{intcmd}、\texttt{car}、\texttt{add}、\texttt{types} 命令；
\item 用串口发送 \texttt{intcmd 7}，观察 LED 动作；
\item 发送 \texttt{car 1.5 2.5 0.8}，观察电机或串口回显；
\item 发送 \texttt{types 12 -3 0x10 3.5 led}，确认动态类型识别正确；
\item 测试超长行、空命令、未知命令不会死机。
\end{enumerate}

\end{anybox}
\begin{anybox}{验收标准}
\begin{itemize}
\item 每条命令有清晰回显与错误提示；
\item 输入非法参数时不会调用业务函数；
\item 长时间连续输入命令不会导致缓冲区溢出或 HardFault。
\end{itemize}

\medskip \noindent\textbf{待补充实验记录}：

\begin{itemize}
\item 单片机与 USB-TTL 接线照片；
\item 串口助手输入与输出截图；
\item 非法输入与边界测试日志。
\end{itemize}
\end{anybox}

\begin{figure}[htbp]
\centering
    % TODO: 在此插入串口接线照片
\caption{串口命令解析实验硬件接线（待补充）} \label{fig:ch9-mcu-hw}
\end{figure}

\begin{figure}[htbp]
\centering
    % TODO: 在此插入串口助手运行截图
\caption{串口命令解析运行日志（待补充）} \label{fig:ch9-mcu-log}
\end{figure}
% ---------------------------------------------------------------------
\begin{problemset}
\item 字符串命令解析的第一步通常是“切词”。请说明 \texttt{空格/制表符} 作为分隔符时，如何支持参数中带空格的字符串。
\item 用哈希表管理命令相比顺序 \texttt{if/strcmp} 有什么好处？哈希冲突时，本章采用了什么解决方案？
\item 命令注册表把“命令名、处理函数、帮助文本”放在一起，对新增命令和命令数量增加分别有什么影响？
\item \texttt{Param\_T} 用一个联合体保存多种参数类型，靠什么字段知道“应该按 int 还是按 float 读取”？
\item 本章第 6 节介绍了两种“无适配器”框架的实现思路：一是“参数描述表 + 通用调度器”，二是“函数指针强制转换（利用 ABI 特性）”。请回答：
\begin{enumerate}
	\item 这两种方案分别试图解决什么问题？各自的实现思路是什么？
	\item 第一种方案（参数描述表 + 通用调度器）在裸机环境中有哪些现实代价？
	\item 第二种方案（函数指针强制转换）依赖什么底层条件？在硬浮点 ABI 或 x86 平台下为什么会失效？
	\item 综合比较，本章最终选择“手写适配器”方案的根本原因是什么？
\end{enumerate}
\end{problemset}